
\section{Architecture Integration}

\subsection{Model Support}

A3 supports any transformer with linear or Conv1D layers:

\begin{table}[H]
\centering
\small
\begin{tabular}{lllll}
\toprule
\textbf{Architecture} & \textbf{Attention Projections} & \textbf{FFN Projections} & \textbf{Conv1D?} & \textbf{Verified} \\
\midrule
DistilBERT-base & \texttt{q\_lin, k\_lin, v\_lin, out\_lin} & \texttt{lin1, lin2} & No & Yes \\
RoBERTa-base & \texttt{query, key, value, output.dense} & \texttt{intermediate.dense, output.dense} & No & Yes \\
GPT-2 & \texttt{c\_attn} (fused), \texttt{c\_proj} & \texttt{c\_fc, c\_proj} & Yes & Yes \\
Qwen2.5 / LLaMA & \texttt{q\_proj, k\_proj, v\_proj, o\_proj} & \texttt{gate\_proj, up\_proj, down\_proj} & No & Yes \\
\bottomrule
\end{tabular}
\caption{Architectures validated with A3 factorized compression and their linear projection types.}
\label{tab:architecture_support}
\end{table}

\subsection{Layer Targeting}

A3 compresses all linear layers by default, with the following protections:

\begin{table}[H]
\centering
\small
\begin{tabular}{lll}
\toprule
\textbf{Layer type} & \textbf{Compressed?} & \textbf{Reason} \\
\midrule
\texttt{classifier}, \texttt{score}, \texttt{lm\_head} & Never & Small and task-critical \\
\texttt{pooler}, \texttt{embedding} & Never & Not factorizable linear maps \\
Attention projections & Yes & Target of interest \\
FFN layers & Yes (optional) & Large parameter share \\
LayerNorm & Never & 1D affine scale/bias vectors \\
\bottomrule
\end{tabular}
\caption{Layer compression rules and preservation criteria enforced by the A3 engine.}
\label{tab:layer_rules}
\end{table}

The \texttt{should\_compress\_linear} helper enforces these rules.

\subsection{Handling Different Architectures}

\subsubsection{DistilBERT (Flat)}

Standard linear layers:

\begin{lstlisting}
distilbert.transformer.layer.0.attention.q_lin  # Linear(768, 768)
distilbert.transformer.layer.0.attention.k_lin  # Linear(768, 768)
distilbert.transformer.layer.0.attention.v_lin  # Linear(768, 768)
distilbert.transformer.layer.0.attention.out_lin # Linear(768, 768)
distilbert.transformer.layer.0.ffn.lin1          # Linear(768, 3072)
distilbert.transformer.layer.0.ffn.lin2          # Linear(3072, 768)
\end{lstlisting}

\textbf{A3 compresses each directly.} No special handling needed.

\subsubsection{RoBERTa (Nested)}

Standard linear layers within the nested \texttt{BertAttention}:

\begin{lstlisting}
roberta.encoder.layer.0.attention.self.query  # Linear(768, 768)
roberta.encoder.layer.0.attention.self.key    # Linear(768, 768)
roberta.encoder.layer.0.attention.self.value  # Linear(768, 768)
roberta.encoder.layer.0.attention.output.dense # Linear(768, 768)
roberta.encoder.layer.0.intermediate.dense    # Linear(768, 3072)
roberta.encoder.layer.0.output.dense          # Linear(3072, 768)
\end{lstlisting}

\textbf{All are standard Linear.} No special handling.

\subsubsection{GPT-2 (Conv1D)}

GPT-2 uses \texttt{Conv1D} from HuggingFace, which stores weights as 
\texttt{(in\_features, out\_features)} instead of \texttt{(out\_features, in\_features)}:

\begin{lstlisting}
transformer.h.0.attn.c_attn  # Conv1D(768, 2304) -- fused Q,K,V
transformer.h.0.attn.c_proj  # Conv1D(768, 768)
transformer.h.0.mlp.c_fc     # Conv1D(768, 3072)
transformer.h.0.mlp.c_proj   # Conv1D(3072, 768)
\end{lstlisting}

\textbf{A3 handles Conv1D} by transposing the weight before SVD:

\begin{lstlisting}
if "conv1d" in cls_name or hasattr(module, "nf"):
    W = raw_w.t()  # (in, out) $\to$ (out, in)
\end{lstlisting}

The factorized layer replaces the Conv1D. Because the factorized layer 
implements a standard linear operation, the \texttt{nf}/\texttt{nx} attributes are preserved 
for compatibility.

\subsubsection{Qwen3.5 (GQA + RoPE)}

Standard PyTorch Linear layers with GQA structure:

\begin{lstlisting}
model.layers.0.self_attn.q_proj  # Linear(1024, 2048) -- H * d_h
model.layers.0.self_attn.k_proj  # Linear(1024, 512)  -- H_kv * d_h
model.layers.0.self_attn.v_proj  # Linear(1024, 512)
model.layers.0.self_attn.o_proj  # Linear(2048, 1024)
\end{lstlisting}

\textbf{A3 compresses each directly.} GQA-specific handling is unnecessary because 
the factorization is per-projection.

\subsection{Wrapper Compatibility}

A3 supports models wrapped with LoRA-style adapters (\texttt{LoRALinear}, etc.):

\begin{lstlisting}
class LoRALinear(nn.Module):
    def __init__(...):
        self.base_layer = nn.Linear(...)  # $\leftarrow$ A3 targets this
        self.lora_A = nn.Parameter(...)
        self.lora_B = nn.Parameter(...)
\end{lstlisting}

\textbf{Injection Procedure.} The injection procedure isolates the underlying dense representation by extracting \texttt{target\_mod = getattr(module, "base\_layer", module)}, factorizing this core weight tensor, replacing \texttt{module.base\_layer} with the factorized linear module, and strictly preserving the outer parameter tensors (\texttt{module.lora\_A} and \texttt{module.lora\_B}). If the target module was previously unadapted, A3 optionally instantiates internal LoRA parameters within the factorized module whenever \texttt{lora\_rank > 0} to maintain downstream fine-tuning capability.

\subsection{Injection Algorithm}

\textbf{Algorithm 2:} Inject A3 into a model
\begin{lstlisting}
Input:  model M, config c, calibration data D
Output: compressed model M'

1: hooks $\leftarrow$ []
2: for   in linear_modules(M) do
3:     hooks.append(register_forward_hook( , accumulate_F_C))
4: end for
5: run_calibration(M, D)
6: for h in hooks: h.remove()
7:
8: F_map $\leftarrow$ finalized_statistics()
9: for   in linear_modules(M) do
10:     if not should_compress( ): continue
11:     U, V, err $\leftarrow$ weighted_svd( .weight, F_map[ ], r, $\tau$)
12:     if err $\le$ $\tau$:
13:         replace( , A3FactorizedLinear(U, V))
14:     else:
15:         mark_fallback( )
16: end for
17:
18: if recalibrate_ln:
19:     $\sigma$_target $\leftarrow$ record_LN_before(M, D)    pre-computed
20:     $\sigma$_current $\leftarrow$ record_LN_after(M', D)
21:     for $\gamma$ in layer_norms(M') do
22:         $\gamma$ $\leftarrow$ $\gamma$ $\cdot$ $\sigma$_current / $\sigma$_target
23:     end for
24: end if
25: return M'
\end{lstlisting}

\subsection{The `A3FactorizedLinear` Module}

The factorized layer implements:

\begin{lstlisting}
class A3FactorizedLinear(nn.Module):
    def __init__(self, in_features, out_features, rank, ...):
        self.V = nn.Parameter(torch.empty(rank, in_features))  # (r, d_in)
        self.U = nn.Parameter(torch.empty(out_features, rank)) # (d_out, r)
        if bias:
            self.bias = nn.Parameter(torch.empty(out_features))
        # Optional LoRA
        if lora_rank > 0:
            self.lora_A = nn.Parameter(...)
            self.lora_B = nn.Parameter(...)
    
    def forward(self, x):
        latent = F.linear(x, self.V)      # (B, L, r)
        out = F.linear(latent, self.U)    # (B, L, d_out)
        if self.lora_rank > 0:
            out += F.linear(F.linear(x, self.lora_A), self.lora_B) * self.lora_scaling
        if self.bias is not None:
            out += self.bias
        return out
\end{lstlisting}

\textbf{Parameter count:} \texttt{r \textit{ (d\_in + d\_out)} vs \texttt{d\_in } d\_out} for the original.

\textbf{Injection procedure:}
The injection procedure safely decouples base weights from adapter representations by extracting \texttt{target\_mod = getattr(module, "base\_layer", module)}, factorizing the underlying dense projection, replacing \texttt{module.base\_layer} with the factorized low-rank module, and strictly preserving the outer LoRA parameter tensors (\texttt{module.lora\_A} and \texttt{module.lora\_B}). If the target module was previously unadapted, A3 optionally instantiates internal LoRA parameters within the factorized module whenever \texttt{lora\_rank > 0} to maintain downstream fine-tuning capability.

\subsection{Injection Algorithm}

\textbf{Algorithm 2:} Inject A3 into a model
\begin{lstlisting}
Input:  model M, config c, calibration data D
Output: compressed model M'

1: hooks $\leftarrow$ []
2: for   in linear_modules(M) do
3:     hooks.append(register_forward_hook( , accumulate_F_C))
4: end for
5: run_calibration(M, D)
6: for h in hooks: h.remove()
7:
8: F_map $\leftarrow$ finalized_statistics()
9: for   in linear_modules(M) do
10:     if not should_compress( ): continue
11:     U, V, err $\leftarrow$ weighted_svd( .weight, F_map[ ], r, $\tau$)
12:     if err $\le$ $\tau$: replace_module(M,  , U, V)
13: end for
14: return M'
\end{lstlisting}

\subsection{Causality and Streaming}

For autoregressive decoder models such as GPT-2, A3's calibration protocol operates in strict compliance with causal attention constraints. Forward hooks intercept intermediate activations dynamically as tokens propagate, ensuring zero future-token leakage because activations are collected exclusively during standard lower-triangular causal evaluation. We empirically verified that the directional gradient of output tokens at position $i$ with respect to input tokens at position $j > i$ remains identically zero both before and after low-rank module replacement.

\subsection{Failure Modes and Fallbacks}

\begin{table}[H]
\centering
\small
\begin{tabular}{lll}
\toprule
\textbf{Potential Failure} & \textbf{Detection Mechanism} & \textbf{Defensive Response} \\
\midrule
Reconstruction error $> \tau$ & Frobenius norm check & Fallback to original dense layer \\
Weight shape mismatch & Assertion in \texttt{from\_linear} & Skip layer and retain dense weights \\
SVD convergence failure & \texttt{torch.linalg.LinAlgError} & Fallback to unweighted CPU SVD \\
No calibration data & \texttt{F\_map} empty & Unweighted CPU FP32 SVD \\
CUDA OOM during SVD & Proactive CPU computation & Execute entirely on CPU in FP32 \\
NaN in reconstruction & \texttt{torch.isnan} check & Abort factorization and fallback \\
\bottomrule
\end{tabular}
\caption{Defensive exception-handling and fallback behaviors implemented within the A3 engine.}
\label{tab:failure_responses}
\end{table}

\subsection{Registry Integration}

\begin{lstlisting}
@COMPRESSOR_REGISTRY.register("a3")
@COMPRESSOR_REGISTRY.register("a3_compressor")
def a3_compressor_factory(model, comp_config, ...):
    ...
\end{lstlisting}

\subsection{Compatibility with the LoRA Recovery Pipeline}

A3 integrates seamlessly with the framework's post-compression adaptation pipeline. The workflow operates as a four-stage sequential pipeline: first, A3 factorizes the target linear projections of the base model subject to safety fallback gates; second, low-rank LoRA adapter tensors are attached directly to the frozen factorized operators; third, recovery fine-tuning optimizes exclusively the parameter-efficient adapter parameters across a constrained budget of 50 steps; finally, the adapted model undergoes downstream benchmark evaluation. This decoupled design preserves parameter efficiency while matching the standard execution protocol deployed across our benchmark suite.